\begin{frame}
	\frametitle{Matematika 2.del}
\end{frame}

\begin{frame}{Zaporedja, vrste in limite}
	\begin{enumerate}
		\item 
		Naj bo $\sum_{n=1}^\infty a_n $ absolutno konvergentna vrsta in $a_n \ne -1$.
		Dokaži, da je tudi vrsta $\sum_{n=1}^\infty \frac{a_n}{1+a_n}$
		absolutno konvergentna.

		\item
		Izračunaj limito
		\[
			\lim_{x \to \infty} (\sin{\sqrt{x+1}}-\sin{\sqrt{x}})
        \]
		\item
		Za dani zaporedji preveri, ali sta konvergentni.
		\begin{columns}
			\begin{column}{0.5\textwidth}
				\begin{equation*}
					a_n=\underbrace{\sqrt{2+\sqrt{2+\dots+\sqrt{2}}}}_{\text{\normalsize{$n$}-korenov}} 
				\end{equation*}
			\end{column}
			\begin{column}{0.5\textwidth}
				\begin{equation*}
					b_n = \underbrace{\sin(\sin(\dots(\sin 1)\dots))}_{\text{\normalsize{$n$}-sinusov}}
				\end{equation*}
			\end{column}
		\end{columns}
		
	\end{enumerate}
\end{frame}

\begin{frame}{Algebra}
	\begin{enumerate}
		\item
		Vektorja $\vec{c}=\vec{a}+2\vec{b}$ in $\vec{d}=\vec{a}-\vec{b}$
		sta pravokotna in imata dolžino 1. Določi kot med vektorjema $\vec{a}$ in $\vec{b}$.
		\item 
		Izračunaj
		\begin{equation*}
			\left(
				\begin{matrix}
					1 & 2 & 3 & 4 & 5 & 6 \\
					4 & 5 & 2 & 6 & 3 & 1 
				\end{matrix}
			\right)^{2000}
		\end{equation*}
		
	\end{enumerate}
\end{frame}

\begin{frame}{Velika determinanta}
	Izračunaj naslednjo determinanto $2n \times 2n$, ki ima na neoznačenih mestih ničle.
	\begin{equation*}
		\left|
			\begin{matrix}
				1 & & & & 1 & & & & \\
				  & 2 & & & 2 & & & & \\
				  &  & \ddots & & \vdots & & & & \\
				  &  & & $n$-1 & $n$-1 & & & & \\
				 1 & 2 & \cdots& $n$-1 & $n$ & $n$+1 & $n$+2 & \cdots & 2$n$ \\
				 &  & &  & $n$+1 & $n$+1& & & \\
				 &  & &  & $n$+1 & & $n$+2& & \\
				 &  &  & & \vdots & & & \ddots& \\
				 &  &  & & 2$n$ & & & & 2$n$\\
				 

			\end{matrix}
		\right|
	\end{equation*}
\end{frame}

\begin{frame}{Grupe}
	
	Naj bo \[
G = \{ z \in \mathbb{C} \mid z = 2k (\cos(m\pi \sqrt{2}) + i\sin(m\pi \sqrt{2})), \; k,m \in \mathbb{Z} \}
\]
\[
H = \{ (x,y) \in \mathbb{R}^2 \mid x, y \in \mathbb{Z} \}
\]
	
	\begin{enumerate}
		\item
			Pokaži, da je $G$ podgrupa v grupi $(\mathbb{C} \setminus \{0\}, \cdot)$
			neničelnih kompleksnih števil za običajno množenje.
		\item
			Pokaži, da je $H$ podgrupa v aditivni grupi $(\mathbb{R}^2, \cdot)$
			ravninskih vektorjev za običajno seštevanje po komponentah.
		\item
			Pokaži, da je preslikava $f:H\to G$, podana s pravilom
			\[
(x, y) \mapsto 2x \left( \cos(y\pi \sqrt{2}) + i \sin(y\pi \sqrt{2}) \right)
\]
			izomorfizem grup $G$ in $H$.
	\end{enumerate}
\end{frame}
